% Generated by `taskgen course examples` from number_theory.miller_rabin (seed 7).
% Do not edit: change examples.yaml and regenerate.
\begin{example}[тест Миллера--Рабина]\label{ex:miller-rabin-prime}
С помощью теста Миллера-Рабина определить, является ли простым число $N = 193$. Требуется дать ответ с вероятностью ошибки не~более~$\frac{1}{16}$.

\par\noindent\textit{Решение.} $N - 1 = 192 = 2^{6} \cdot 3$. Для основания $a$ вычисляем $x_0 = a^{3}$, $x_i = x_{i-1}^2$; $a$ --- не свидетель, если $x_0 = 1$ или $x_i = N - 1 = 192$ при некотором $i < 6$. Для составного $N$ не-свидетелей не более $1/4$, поэтому нужно два раунда ($(1/4)^2 = 1/16$).\par\noindent $a = 2$: $x_i$: $8, 64, 43, 112, 192$ --- раунд пройден.\par\noindent $a = 3$: $x_i$: $27, 150, 112, 192$ --- раунд пройден.\par\noindent Пройдено 2 раунда, вероятность ошибки не более $\frac{1}{16}$: $N = 193$ --- простое.
\end{example}
